% Part: turing-machines
% Chapter: machines-computations
% Section: halting-states

\documentclass[../../../include/open-logic-section]{subfiles}

\begin{document}

\olfileid{tur}{mac}{hal}
\olsection{حالت‌های توقف}

\begin{explain}
گرچه ماشین‌هایمان را چنان تعریف کرده‌ایم که تنها هنگام نبود دستور
قابل اجرا متوقف شوند، نمایش‌های رایج ماشین‌های تورینگ یک \emph{حالت
توقفِ} اختصاصی~$h$ دارند که $h\in Q$.

اندیشهٔ پشت حالت توقف ساده است: هنگامی که کار ماشین پایان یافته است
(آمادهٔ پذیرفتن ورودی است یا نوشتن خروجی را تمام کرده)، به حالت~$h$
می‌رود و در آن متوقف می‌شود. بعضی ماشین‌ها دو حالت توقف دارند: یکی
برای پذیرفتن ورودی و دیگری برای رد آن.
\end{explain}

\begin{ex}\emph{حالت‌های توقف}.
برای روشن‌کردن این مفهوم، با تغییری در ماشین زوج آغاز کنیم. به‌جای
آنکه ماشین در صورت زوج‌بودن ورودی در حالت~$q_0$ متوقف شود، می‌توانیم
دستوری بیفزاییم که ماشین را به حالت توقف بفرستد.
\[
\begin{tikzpicture}[->,>=stealth',shorten >=1pt,auto,node distance=2.8cm,
                    semithick]
  \tikzstyle{every state}=[fill=none,draw=black,text=black]

  \node[initial, state]         (A)                     {$q_0$};
  \node[state]         (B) [right of=A] {$q_1$};
  \node[state]         (C) [below of=A] {$h$};

  \path (A) edge [bend left] node {\TMtrans{\TMstroke}{\TMstroke}{R}} (B)
        edge node {\TMtrans{\TMblank}{\TMblank}{N}} (C)
        (B) edge [loop above] node {\TMtrans{\TMblank}{\TMblank}{R}} (B)
            edge [bend left] node {\TMtrans{\TMstroke}{\TMstroke}{R}} (A);
\end{tikzpicture}
\]

مثال را بیشتر گسترش دهیم. وقتی ماشین تشخیص می‌دهد ورودی فرد است،
هرگز متوقف نمی‌شود. می‌توانیم با جایگزین‌کردن دستور حلقه با دستوری
برای رفتن به حالت رد~$r$، ماشین را چنان تغییر دهیم که \emph{حالت رد}
داشته باشد.
\[
\begin{tikzpicture}[->,>=stealth',shorten >=1pt,auto,node distance=2.8cm,
                    semithick]
  \tikzstyle{every state}=[fill=none,draw=black,text=black]

  \node[initial,state]         (A)                     {$q_0$};
  \node[state]         (B) [right of=A] {$q_1$};
  \node[state]         (C) [below of=A] {$h$};
  \node[state]         (D) [below of=B] {$r$};

  \path (A) edge [bend left] node {\TMtrans{\TMstroke}{\TMstroke}{R}} (B)
        edge node {\TMtrans{\TMblank}{\TMblank}{N}} (C)
        (B) edge node {\TMtrans{\TMblank}{\TMblank}{N}} (D)
            edge [bend left] node {\TMtrans{\TMstroke}{\TMstroke}{R}} (A);
\end{tikzpicture}
\]
\end{ex}

\begin{explain}
افزودن حالت توقف اختصاصی در مواردی مانند این سودمند است، زیرا زمان
پذیرفتن یا ردکردن بعضی ورودی‌ها را صریح می‌کند. بااین‌حال، توجه به
این نکته مهم است که افزودن حالت توقف اختصاصی هیچ توان محاسباتی‌ای
نمی‌افزاید. به همین ترتیب، مفهوم کم‌تر رسمی توقف نیز مزایای خودش را
دارد. تعریف توقفی که تا اینجا در این فصل به کار برده‌ایم برهان
\emph{مسئلهٔ توقف} را شهودی و نمایش آن را آسان می‌کند. به همین دلیل
با تعریف اصلی خود ادامه می‌دهیم.
\end{explain}

\end{document}
